\section*{Comparative Reception and Theological Dossier}

\begin{studybox}{Method}
The Old Testament and Gospel are officially correlated; Romans is
semi-continuous; the Week XVIII Missal formulary is shared by all three cycles.
The witnesses below illuminate appointed units without being made authors of
the complete Year A packet. Historical setting, patristic reception,
sacramental resonance, and present application remain distinct.
\end{studybox}

\subsection*{Isaiah's market without a price}

Isaiah~55:1--3 addresses thirst, hunger, failed expenditure, hearing, life, and
an everlasting covenant. The sequence of imperatives matters. The hearer is
not offered an inexpensive commodity but summoned away from spending on what
does not satisfy and toward attentive reception of God's word. The Davidic
covenant in verse~3 prevents the invitation from becoming a free-floating
maxim about personal fulfillment. It is restoration speech addressed to a
people whose future depends upon divine fidelity.

Ancient Christian reception commonly extends the water, wine, milk, and bread
images toward baptism, instruction, wisdom, and the Eucharist. Those readings
are strongest when they preserve Isaiah's first claim: God gives life and
renews covenant. They become weaker when every image is forced to signify one
sacrament in one-to-one fashion. Augustine's theology of gratuitous grace is a
near theological analogue rather than a direct commentary on the appointed
cut. The text's own economic verbs also retain moral force. ``Without money''
does not praise exploitation or deny the cost of human labor; it denies that
divine covenant life can be purchased.

Thoreau's explicit reuse in \emph{Walden} demonstrates how far the question can
travel. He redirects Isaiah's language toward nineteenth-century American
consumption and lives spent acquiring unnecessary things. That is genuine
verbal reception in a changed register, but Thoreau's economic program is not
therefore the prophet's historical message.

\subsection*{The open hand and the social field of praise}

The appointed Psalm~145 verses can sound like a compact theology of universal
provision. The complete psalm is more socially specific: the Lord executes
justice for the oppressed, feeds the hungry, frees prisoners, raises the bowed
down, loves the just, protects strangers, and sustains widow and orphan.
Creation's expectation and food in due season therefore stand inside praise of
a reign that includes vulnerable persons and frustrates wicked ways.

Augustine's exposition of Psalm~144 receives the expectant creatures as
dependent upon the one giver. That dependence is not passivity. In Matthew,
Christ refuses to make the disciples spectators: their proposal dismisses the
crowd, but his command makes their acknowledged insufficiency the point of
departure for service. Read together, psalm and Gospel distinguish divine
source from creaturely mediation. God opens the hand; disciples receive,
distribute, and gather.

This relation also sets a limit on prosperity readings. The psalm is praise,
not a guarantee that every faithful person will possess adequate food on
demand. The Gospel's abundance cannot be used to deny famine, economic
injustice, or the obligation to organize ordinary material assistance.

\subsection*{Matthew's two banquets and the training of disciples}

The feeding follows Matthew's account of John the Baptist's death. Herod's
banquet is enclosed, elite, politically anxious, and ends with a murdered
prophet's head given on a dish. Jesus withdraws, receives the pursuing crowd
with compassion, heals the sick, seats the people on grass, blesses and breaks
what is present, and feeds all. The narrative contrast is not stated as a
formal allegory, yet the adjacent scenes make two kinds of kingship and table
visible.

Chrysostom's Homily~49 follows the disciples' education closely. Their
insufficiency is disclosed rather than mocked; Christ could feed without them,
yet makes them ministers and has fragments gathered. His reading supports
providence, discipline, service, and avoidance of waste. It does not authorize
the claim that the disciples secretly possessed enough, or that the miracle
was only an exercise in induced sharing. Nor does Matthew identify this meal
with the Eucharist. The verbs of taking, blessing, breaking, and giving create
an intelligible Eucharistic reception horizon, while the narrated action
remains a wilderness feeding of bodily hungry people.

The twelve baskets have generated Israel-and-apostles readings, but the text
itself first establishes abundance after universal satisfaction. Symbolic
reception should not erase the practical command to gather what remains.

\subsection*{Romans: inseparability amid—not exemption from—affliction}

Romans~8:35, 37--39 is discontinuous in proclamation: verse~36's scriptural
citation about being killed all day is omitted. The wider argument must remain
in view. Tribulation, distress, persecution, famine, nakedness, peril, and
sword are not imaginary; the claim is that none can separate those in Christ
from divine love. ``More than conquerors'' therefore does not mean invulnerable,
socially dominant, or exempt from grief.

Chrysostom emphasizes Paul's catalogue as proof of love tested in suffering,
not a celebration of pain. Aquinas distinguishes the threatening forces while
locating security in God's initiative. The passage cannot require a person to
remain in abuse, refuse medicine, conceal crime, or avoid lawful protection.
Its hope concerns communion with God through real affliction, not the moral
neutralization of perpetrators or dangerous conditions.

Arnold's \emph{Dover Beach} creates a revealing changed-register inversion.
Its accumulating ``neither/nor'' world lacks certitude, peace, and help for
pain, so fidelity contracts toward two lovers. Romans' catalogue also moves
through negations, but reaches the expansive prior love of God in Christ. The
formal echo is illuminating even though conscious borrowing has not been
established.

\subsection*{One liturgy, several kinds of nourishment}

The packet should not be reduced to ``bread'' as a single undifferentiated
symbol. Isaiah joins nourishment to hearing and covenant; the psalm praises
the creator's timely provision; Matthew narrates bodily feeding through
compassionate action; the acclamation says life exceeds bread alone; Communion
A receives wilderness food through Wisdom; Communion B names Christ as bread
of life. The distinctions are the theology.

The shared Week XVIII prayers add another grammar. Created persons require
restoration and preservation; gifts require sanctification; worshippers ask to
become an enduring offering; communicants receive a heavenly gift and seek
continuing protection. The two Communion antiphons remain alternatives. If
Wisdom is used, manna-reception stands at the ending; if John is used, Christ's
self-identification stands there. A guide may explain both branches but may
not combine them into an enacted text without evidence of the actual
selection.

\begin{threestable}{Convergence}{Deepened claim}{Guardrail}
Invitation, expectation, compassion & Divine giving precedes and enables human response. & Gratuitous gift does not erase labor, justice, or organized care.\\
Loaves, word, and heavenly bread & Bodily and spiritual hungers are held together without identity. & Neither spiritualize material need nor reduce salvation to present sufficiency.\\
Distribution and self-offering & Disciples and worshippers become agents within a gift they did not originate. & Service confers no divine ownership or permission to coerce.\\
Affliction and inseparable love & Hope persists inside named dangers. & Never use Romans to demand endurance of preventable harm.\\
\end{threestable}
